\section{the model description; notations}
\label{sec:the-model-description-notations}
we consider the system consisting of the free alpha particle and one atom, described as a 1D harmonic oscillator, located at $x=a_1$.

the system Hamiltonian is
\begin{align}
  H=\frac{P_X^2}{2M}+\frac{P_x^2}{2m}+\frac{1}{2}m\omega_0^2(x-a_1)^2+V_0e^{-\frac{(x-X)^2}{2\sigma^2}}=H_0+H_I
\end{align}
where the unperturbed Hamiltonian is given by
\begin{align}
  H_0=\frac{P_X^2}{2M}+\frac{P_x^2}{2m}+\frac{1}{2}m\omega_0^2(x-a_1)^2
\end{align}
and the perturbation is given by
\begin{align}
  H_I=V_0e^{-\frac{(x-X)^2}{2\sigma^2}}
\end{align}
The energy eigenvectors of $H_0$ are
\begin{align}
  H_0|n(a_1),P\rangle = (\hbar\omega_0(n+1/2)+\frac{P^2}{2M})|n(a_1),P\rangle
\end{align}
in the position representations they are given by
\begin{align}
  \langle x,X|n(a_1),P\rangle = u_n(x-a_1)\frac{1}{\sqrt{2\pi\hbar}}e^{\frac{i}{\hbar}PX}
\end{align}
the initial state of the system in the position representation is given by
\begin{align}
  \langle x,X|\psi_i\rangle = u_0(x-a_1) \psi_\alpha(X)
\end{align}
where $\psi_\alpha(X)$ is the bimodal  wave packet of the alpha particle with initial momentum $P_0$ and initial width $\sigma_0$ described in Sect. \ref{sect_bimodal}
\begin{align}
  \psi_\alpha(X) = \frac{1}{\sqrt{\sigma_\alpha\sqrt{\pi}}}\sqrt{\frac{2}{1+e^{-\frac{P_0^2\sigma_\alpha^2}{\hbar^2}}}}e^{-\frac{X^2}{2\sigma_\alpha^2}}\cos\left(\frac{P_0 X}{\hbar}\right)\label{psi(X_0)-tdse_tdpt}
\end{align}
the energy eigenfunctions expansion of $\psi_\alpha$ is
\begin{align}
  \psi_\alpha(X)=\frac1{\sqrt{2\pi\hbar}}\int\limits_{-\infty}^{\infty}dP C(P)e^{\frac{iPX}{\hbar}}
\end{align}
where
\begin{align}
  C(P) = \langle P|\psi_\alpha\rangle = \frac{1}{\sqrt{2\pi\hbar}}\int\limits_{-\infty}^{\infty}dX e^{-\frac{i}{\hbar}PX}\psi_\alpha(X)
\end{align}
Substituting the explicit form of $\psi_\alpha(X)$ yields:
\begin{align}
  C(P) = \frac{1}{\sqrt{2\pi\hbar}}\frac{1}{\sqrt{\sigma_\alpha\sqrt{\pi}}}\sqrt{\frac{2}{1+e^{-\frac{P_0^2\sigma_\alpha^2}{\hbar^2}}}}\int\limits_{-\infty}^{\infty}dX e^{-\frac{i}{\hbar}PX}e^{-\frac{X^2}{2\sigma_\alpha^2}}\cos\left(\frac{P_0 X}{\hbar}\right)
\end{align}
Using $\cos(u) = \frac{1}{2}(e^{iu} + e^{-iu})$, we can evaluate the integral:
\begin{align}
  \int\limits_{-\infty}^{\infty}dX e^{-\frac{i}{\hbar}PX}e^{-\frac{X^2}{2\sigma_\alpha^2}}\cos\left(\frac{P_0 X}{\hbar}\right) & = \frac{1}{2}\int\limits_{-\infty}^{\infty}dX e^{-\frac{X^2}{2\sigma_\alpha^2}}\left[ e^{-\frac{i}{\hbar}(P-P_0)X} + e^{-\frac{i}{\hbar}(P+P_0)X} \right] \nonumber \\
                                                                                                                               & = \frac{1}{2}\sqrt{2\pi\sigma_\alpha^2}\left[ e^{-\frac{(P-P_0)^2\sigma_\alpha^2}{2\hbar^2}} + e^{-\frac{(P+P_0)^2\sigma_\alpha^2}{2\hbar^2}} \right]
\end{align}
Substituting this back into the expression for $C(P)$ and simplifying the coefficients, we obtain the final result:
\begin{align}
  C(P) = \sqrt{\frac{\sigma_\alpha}{2\hbar\sqrt{\pi}\left(1+e^{-\frac{P_0^2\sigma_\alpha^2}{\hbar^2}}\right)}}\left[ e^{-\frac{(P-P_0)^2\sigma_\alpha^2}{2\hbar^2}} + e^{-\frac{(P+P_0)^2\sigma_\alpha^2}{2\hbar^2}} \right]\label{C(P)-one-atom}
\end{align}

Note that for $P_0 \gg \hbar/\sigma_\alpha$, the second term in the brackets is negligible, and we can approximate $C(P)$ as:
\begin{align}
  C(P) \approx \sqrt{\frac{\sigma_\alpha}{2\hbar\sqrt{\pi}}}\left[ e^{-\frac{(P-P_0)^2\sigma_\alpha^2}{2\hbar^2}}+ e^{-\frac{(P+P_0)^2\sigma_\alpha^2}{2\hbar^2}} \right]\label{C(P)-one-atom-approx}
\end{align}

